\documentclass[UTF8, 12pt, fontset=fandol, oneside]{ctexart}
\usepackage{researchnotes}

\title{第六章\quad 特征值}
\author{}
\date{}

\begin{document}

\maketitle

\section*{§6.1\quad 特征值和特征向量}

在这一章及下一章中，我们主要研究有限维线性空间上的线性变换. 我们特别关心这样一个问题：给定线性空间 $V$ 上的线性变换 $\varphi$，能否找到 $V$ 的一组基，使线性变换 $\varphi$ 在这组基下的表示矩阵具有特别简单的形状？比如，若我们能找到 $V$ 的一组基 $\{e_1,e_2,\cdots,e_n\}$，使线性变换 $\varphi$ 在这组基下的表示矩阵为对角阵：
\[
\left(\begin{array}{cccc}
a_1 & & & \\
& a_2 & & \\
& & \ddots & \\
& & & a_n
\end{array}\right).
\]
这时，若 $\alpha=k_1e_1+k_2e_2+\cdots+k_ne_n$，则
\[
\varphi(\alpha)=a_1k_1e_1+a_2k_2e_2+\cdots+a_nk_ne_n.
\]
线性变换 $\varphi$ 的表达式非常简单，线性变换 $\varphi$ 的许多性质也变得一目了然. 例如，若 $a_1,a_2,\cdots,a_r$ 不为零，而 $a_{r+1}=\cdots=a_n=0$，则 $\varphi$ 的秩为 $r$，且 $\operatorname{Im}\varphi$ 就是由 $\{e_1,e_2,\cdots,e_r\}$ 生成的子空间，而 $\operatorname{Ker}\varphi$ 则是由 $\{e_{r+1},\cdots,e_n\}$ 生成的子空间.

由第四章我们已经知道，一个线性变换在不同基下的表示矩阵是相似的. 因此用矩阵的语言重述上面提到的问题就是：能否找到一类特别简单的矩阵，使任一矩阵都与这类矩阵中的某一个相似？比如，我们可以问：是否所有的矩阵都相似于对角阵？若不然，哪一类矩阵可以相似于对角阵？

由第四章第 5 节知道，若线性空间 $V$ 可分解为
\begin{equation}
V=V_1\oplus V_2\oplus\cdots\oplus V_m,
\tag{6.1.1}
\end{equation}
其中每个 $V_i$ 都是线性变换 $\varphi$ 的不变子空间，那么 $\varphi$ 可以表示为分块对角阵. 我们当然希望 (6.1.1) 式中的 $V_i$ 越小越好. 最小的非零子空间是一维子空间. 若 $V_i$ 是一维子空间，$x$ 是其中的任一非零向量，$\varphi$ 在 $V_i$ 上的作用相当于一个数乘，于是存在 $\lambda_0\in\mathbb{K}$，使
\[
\varphi(x)=\lambda_0x.
\]

\noindent\textbf{定义 6.1.1}\quad 设 $\varphi$ 是数域 $\mathbb{K}$ 上线性空间 $V$ 上的线性变换，若 $\lambda_0\in\mathbb{K}$，$x\in V$ 且 $x\neq 0$，使
\begin{equation}
\varphi(x)=\lambda_0x,
\tag{6.1.2}
\end{equation}
则称 $\lambda_0$ 是线性变换 $\varphi$ 的一个特征值，向量 $x$ 称为 $\varphi$ 关于特征值 $\lambda_0$ 的特征向量.

由 (6.1.2) 式我们可以看出，$\varphi$ 关于特征值 $\lambda_0$ 的全体特征向量再加上零向量构成 $V$ 的一个子空间. 事实上，若向量 $x,y$ 是关于特征值 $\lambda_0$ 的特征向量，则
\[
\varphi(x+y)=\varphi(x)+\varphi(y)=\lambda_0x+\lambda_0y=\lambda_0(x+y),
\]
\[
\varphi(cx)=c\varphi(x)=c\lambda_0x=\lambda_0(cx).
\]
因此 $\varphi$ 的关于特征值 $\lambda_0$ 的全体特征向量加上零向量构成 $V$ 的子空间，记为 $V_{\lambda_0}$，称为 $\varphi$ 的关于特征值 $\lambda_0$ 的特征子空间. 显然 $V_{\lambda_0}$ 是 $\varphi$ 的不变子空间.

现在设 $\varphi$ 在某组基下的表示矩阵为 $A$，向量 $x$ 在这组基下可表示为一个列向量 $\alpha$. 这时 (6.1.2) 式等价于
\begin{equation}
A\alpha=\lambda_0\alpha,
\tag{6.1.3}
\end{equation}
(6.1.3) 式也等价于
\begin{equation}
(\lambda_0I_n-A)\alpha=0.
\tag{6.1.4}
\end{equation}

\noindent\textbf{定义 6.1.2}\quad 设 $A$ 是数域 $\mathbb{K}$ 上的 $n$ 阶方阵，若存在 $\lambda_0\in\mathbb{K}$ 及 $n$ 维非零列向量 $\alpha$，使 (6.1.3) 式成立，则称 $\lambda_0$ 为矩阵 $A$ 的一个特征值，$\alpha$ 为 $A$ 关于特征值 $\lambda_0$ 的特征向量. 齐次线性方程组 $(\lambda_0I_n-A)x=0$ 的解空间 $V_{\lambda_0}$ 称为 $A$ 关于特征值 $\lambda_0$ 的特征子空间.

我们已经定义了线性变换与矩阵的特征值，现在的问题是如何来求一个线性变换或一个矩阵的特征值？从 (6.1.4) 式可以看出，要使 $\alpha$ 非零，必须 $|\lambda_0I_n-A|=0$. 反过来，若 $\lambda_0\in\mathbb{K}$ 且 $|\lambda_0I_n-A|=0$，则 (6.1.4) 式有非零解 $\alpha$. 因此寻找矩阵 $A$ 的特征值等价于寻找行列式 $|\lambda I_n-A|=0$ 时 $\lambda$ 的值. 设 $A=(a_{ij})$，则
\begin{equation}
|\lambda I_n-A|=
\left|\begin{array}{cccc}
\lambda-a_{11} & -a_{12} & \cdots & -a_{1n}\\
-a_{21} & \lambda-a_{22} & \cdots & -a_{2n}\\
\vdots & \vdots & & \vdots\\
-a_{n1} & -a_{n2} & \cdots & \lambda-a_{nn}
\end{array}\right|
\tag{6.1.5}
\end{equation}
是一个以 $\lambda$ 为未知数的 $n$ 次首一多项式.

\noindent\textbf{定义 6.1.3}\quad 设 $A$ 是 $n$ 阶方阵，称 $|\lambda I_n-A|$ 为 $A$ 的特征多项式.

由上面的讨论可得矩阵 $A$ 的特征值就是它的特征多项式的根. 读者会提出这样的问题：既然同一个线性变换在不同基下的表示矩阵是相似的，那么相似矩阵是否有相同的特征值？回答是肯定的，这就是下面的定理.

\noindent\textbf{定理 6.1.1}\quad 若 $B$ 与 $A$ 相似，则 $B$ 与 $A$ 具有相同的特征多项式，从而具有相同的特征值 (计重数).

\noindent\textbf{证明}\quad 设 $B=P^{-1}AP$，其中 $P$ 是可逆阵，则
\[
|\lambda I_n-B|=|P^{-1}(\lambda I_n-A)P|=|P^{-1}||\lambda I_n-A||P|=|\lambda I_n-A|.
\]
因此相似矩阵必有相同的特征多项式，从而必有相同的特征值 (计重数). $\square$

\noindent\textbf{定义 6.1.4}\quad 设 $\varphi$ 是线性空间 $V$ 上的线性变换，$\varphi$ 在 $V$ 的某组基下的表示矩阵为 $A$，由定理 6.1.1 知 $|\lambda I_n-A|$ 与基或表示矩阵的选取无关，称 $|\lambda I_n-A|$ 为 $\varphi$ 的特征多项式，记为 $|\lambda I_V-\varphi|$.

设
\[
\begin{aligned}
|\lambda I_n-A|&=\lambda^n+a_1\lambda^{n-1}+\cdots+a_{n-1}\lambda+a_n\\
&=(\lambda-\lambda_1)(\lambda-\lambda_2)\cdots(\lambda-\lambda_n).
\end{aligned}
\]
由 Vieta 定理知 $\lambda_1+\lambda_2+\cdots+\lambda_n=-a_1$，$\lambda_1\lambda_2\cdots\lambda_n=(-1)^na_n$. 由行列式 (6.1.5) 不难看出 $a_1=-(a_{11}+a_{22}+\cdots+a_{nn})=-\operatorname{tr}(A)$，$a_n=(-1)^n|A|$. 因此 $A$ 的 $n$ 个特征值的和与积分别为
\[
\lambda_1+\lambda_2+\cdots+\lambda_n=\operatorname{tr}(A),
\]
\[
\lambda_1\lambda_2\cdots\lambda_n=|A|.
\]
从上面的分析我们可以得出求一个矩阵的特征值与特征向量的方法：作矩阵 $\lambda I_n-A$ (通常称为 $A$ 的特征矩阵) 并求出特征多项式 $|\lambda I_n-A|$ 的根，这就是 $A$ 的特征值. 将每个特征值代入线性方程组
\[
(\lambda I_n-A)x=0,
\]
求出非零解，就是对应的特征向量.

\noindent\textbf{例 6.1.1}\quad 设 $A$ 是一个上三角阵：
\[
\left(\begin{array}{cccc}
a_{11} & a_{12} & \cdots & a_{1n}\\
0 & a_{22} & \cdots & a_{2n}\\
\vdots & \vdots & & \vdots\\
0 & 0 & \cdots & a_{nn}
\end{array}\right).
\]
求 $A$ 的特征值.

\noindent\textbf{解}\quad $|\lambda I_n-A|$ 是一个上三角行列式，因此
\[
|\lambda I_n-A|=(\lambda-a_{11})(\lambda-a_{22})\cdots(\lambda-a_{nn}),
\]
即 $A$ 的特征值等于 $A$ 主对角线上的元素 $a_{11},a_{22},\cdots,a_{nn}$. 对下三角阵也有类似的结论.

\noindent\textbf{例 6.1.2}\quad 求下列三阶矩阵的特征值与特征向量：
\[
A=\left(\begin{array}{ccc}
3 & 1 & -1\\
2 & 2 & -1\\
2 & 2 & 0
\end{array}\right).
\]

\noindent\textbf{解}\quad $A$ 的特征多项式为
\[
\left|\begin{array}{ccc}
\lambda-3 & -1 & 1\\
-2 & \lambda-2 & 1\\
-2 & -2 & \lambda
\end{array}\right|
=\lambda^3-5\lambda^2+8\lambda-4=(\lambda-1)(\lambda-2)^2,
\]
因此 $A$ 的特征值为 $1,2$. 设特征向量
\[
x=\left(\begin{array}{c}
x_1\\
x_2\\
x_3
\end{array}\right),
\]
将 $\lambda=1$ 代入 $(\lambda I_3-A)x=0$ 得
\[
\left\{\begin{array}{l}
-2x_1-x_2+x_3=0,\\
-2x_1-x_2+x_3=0,\\
-2x_1-2x_2+x_3=0.
\end{array}\right.
\]
这个方程组系数矩阵的秩为 2，故只有一个线性无关的解，可取为
\[
\xi_1=\left(\begin{array}{c}
1\\
0\\
2
\end{array}\right).
\]
同理将 $\lambda=2$ 代入 $(\lambda I_3-A)x=0$，解得 (也只有一个线性无关的解)
\[
\xi_2=\left(\begin{array}{c}
1\\
1\\
2
\end{array}\right).
\]
故关于特征值 1 的特征向量为 $c_1\xi_1$，关于特征值 2 的特征向量为 $c_2\xi_2$，其中 $c_1,c_2$ 为 $\mathbb{K}$ 中任意的非零数.

\noindent\textbf{例 6.1.3}\quad 求下列矩阵的特征值：
\[
A=\left(\begin{array}{cc}
0 & -1\\
1 & 0
\end{array}\right).
\]

\noindent\textbf{解}\quad 因为
\[
\left|\begin{array}{cc}
\lambda & 1\\
-1 & \lambda
\end{array}\right|=\lambda^2+1,
\]
所以 $A$ 的特征值为 $\mathrm{i},-\mathrm{i}$.

例 6.1.3 表明，即使是有理数域上的矩阵，其特征值有可能是虚数. 这就是说，对数域 $\mathbb{K}$ 上的矩阵 (或相应的线性变换)，有可能在 $\mathbb{K}$ 中不存在特征值. 但是对复数域来说，任一 $n$ 阶方阵总存在特征值. 因此在考虑特征值问题时，我们常常放在复数域里讨论.

我们从例 6.1.1 也看到，一个上三角 (或下三角) 阵的特征值都在主对角线上. 如果能把一个矩阵相似地变换到一个上三角阵，那么它的特征值也就一目了然了. 但是，由于一个矩阵的特征值有可能是虚数，故数域 $\mathbb{K}$ 上的矩阵未必能相似于一个上三角阵. 然而复数域 $\mathbb{C}$ 上的矩阵总相似于上三角 (或下三角) 阵.

\noindent\textbf{定理 6.1.2}\quad 任一复方阵必相似于一上三角阵.

\noindent\textbf{证明}\quad 设 $A$ 是 $n$ 阶复方阵，现对 $n$ 用数学归纳法. 当 $n=1$ 时结论显然成立. 假设对 $n-1$ 阶矩阵结论成立，现对 $n$ 阶矩阵 $A$ 来证明. 设 $\lambda_1$ 是 $A$ 的一个特征值，则存在非零列向量 $\alpha_1$，使
\[
A\alpha_1=\lambda_1\alpha_1.
\]
将 $\alpha_1$ 作为 $\mathbb{C}_n$ 的一个基向量，并扩展为 $\mathbb{C}_n$ 的一组基 $\{\alpha_1,\alpha_2,\cdots,\alpha_n\}$. 将这些基向量按照列分块方式拼成矩阵 $P=(\alpha_1,\alpha_2,\cdots,\alpha_n)$，则 $P$ 为 $n$ 阶非异阵，且
\[
\begin{aligned}
AP&=A(\alpha_1,\alpha_2,\cdots,\alpha_n)=(A\alpha_1,A\alpha_2,\cdots,A\alpha_n)\\
&=(\alpha_1,\alpha_2,\cdots,\alpha_n)
\left(\begin{array}{cc}
\lambda_1 & *\\
O & A_1
\end{array}\right),
\end{aligned}
\]
其中 $A_1$ 是一个 $n-1$ 阶方阵. 注意到 $P=(\alpha_1,\alpha_2,\cdots,\alpha_n)$ 非异，上式即为
\[
P^{-1}AP=\left(\begin{array}{cc}
\lambda_1 & *\\
O & A_1
\end{array}\right).
\]
因为 $A_1$ 是一个 $n-1$ 阶方阵，所以由归纳假设可知，存在 $n-1$ 阶非异阵 $Q$，使 $Q^{-1}A_1Q$ 是一个上三角阵. 令
\[
R=\left(\begin{array}{cc}
1 & O\\
O & Q
\end{array}\right),
\]
则 $R$ 是 $n$ 阶非异阵，且
\[
\begin{aligned}
R^{-1}P^{-1}APR
&=
\left(\begin{array}{cc}
1 & O\\
O & Q
\end{array}\right)^{-1}
\left(\begin{array}{cc}
\lambda_1 & *\\
O & A_1
\end{array}\right)
\left(\begin{array}{cc}
1 & O\\
O & Q
\end{array}\right)\\
&=
\left(\begin{array}{cc}
1 & O\\
O & Q^{-1}
\end{array}\right)
\left(\begin{array}{cc}
\lambda_1 & *\\
O & A_1
\end{array}\right)
\left(\begin{array}{cc}
1 & O\\
O & Q
\end{array}\right)\\
&=
\left(\begin{array}{cc}
\lambda_1 & *\\
O & Q^{-1}A_1Q
\end{array}\right).
\end{aligned}
\]
这是一个上三角阵，它与 $A$ 相似. $\square$

\noindent\textbf{注}\quad 虽然一般数域 $\mathbb{K}$ 上的矩阵未必相似于上三角阵，但是从定理 6.1.2 的证明可以看出，若数域 $\mathbb{K}$ 上的 $n$ 阶方阵 $A$ 的特征值全在 $\mathbb{K}$ 中，则存在 $\mathbb{K}$ 上的非异阵 $P$，使 $P^{-1}AP$ 是一个上三角阵.

作为定理 6.1.2 的应用，我们来证明 3 个有用的命题. 首先，若 $A$ 是一个 $n$ 阶矩阵，$f(x)=a_mx^m+a_{m-1}x^{m-1}+\cdots+a_1x+a_0$ 是一个多项式，记
\[
f(A)=a_mA^m+a_{m-1}A^{m-1}+\cdots+a_1A+a_0I_n.
\]
我们来考虑矩阵 $A$ 的特征值与矩阵 $f(A)$ 的特征值之间的关系.

\noindent\textbf{命题 6.1.1}\quad 设 $n$ 阶矩阵 $A$ 的全部特征值为 $\lambda_1,\lambda_2,\cdots,\lambda_n$，$f(x)$ 是一个多项式，则 $f(A)$ 的全部特征值为 $f(\lambda_1),f(\lambda_2),\cdots,f(\lambda_n)$.

\noindent\textbf{证明}\quad 因为任一 $n$ 阶矩阵均复相似于上三角阵，可设
\[
P^{-1}AP=
\left(\begin{array}{cccc}
\lambda_1 & * & \cdots & *\\
0 & \lambda_2 & \cdots & *\\
\vdots & \vdots & & \vdots\\
0 & 0 & \cdots & \lambda_n
\end{array}\right).
\]
因为上三角阵的和、数乘及乘方仍是上三角阵，经计算不难得到
\[
P^{-1}f(A)P=f(P^{-1}AP)=
\left(\begin{array}{cccc}
f(\lambda_1) & * & \cdots & *\\
0 & f(\lambda_2) & \cdots & *\\
\vdots & \vdots & & \vdots\\
0 & 0 & \cdots & f(\lambda_n)
\end{array}\right).
\]
因此 $f(A)$ 的全部特征值为 $f(\lambda_1),f(\lambda_2),\cdots,f(\lambda_n)$. $\square$

\noindent\textbf{命题 6.1.2}\quad 设 $n$ 阶矩阵 $A$ 适合一个多项式 $g(x)$，即 $g(A)=O$，则 $A$ 的任一特征值 $\lambda_0$ 也必适合 $g(x)$，即 $g(\lambda_0)=0$.

\noindent\textbf{证明}\quad 设 $\alpha$ 是 $A$ 关于特征值 $\lambda_0$ 的特征向量，经简单计算得
\[
g(\lambda_0)\alpha=g(A)\alpha=0.
\]
而 $\alpha\neq 0$，因此 $g(\lambda_0)=0$. $\square$

对可逆阵 $A$，其逆阵 $A^{-1}$ 的特征值和 $A$ 的特征值有什么关系呢？下面的命题回答了这个问题.

\noindent\textbf{命题 6.1.3}\quad 设 $n$ 阶矩阵 $A$ 是可逆阵，且 $A$ 的全部特征值为 $\lambda_1,\lambda_2,\cdots,\lambda_n$，则 $A^{-1}$ 的全部特征值为 $\lambda_1^{-1},\lambda_2^{-1},\cdots,\lambda_n^{-1}$.

\noindent\textbf{证明}\quad 首先注意到 $A$ 是可逆阵，$\lambda_1\lambda_2\cdots\lambda_n=|A|\neq 0$，因此每个 $\lambda_i\neq 0$ (事实上，$A$ 可逆的充分必要条件是它的特征值全不为零).

由定理 6.1.2 可设
\[
P^{-1}AP=
\left(\begin{array}{cccc}
\lambda_1 & * & \cdots & *\\
0 & \lambda_2 & \cdots & *\\
\vdots & \vdots & & \vdots\\
0 & 0 & \cdots & \lambda_n
\end{array}\right).
\]
因为上三角阵的逆阵仍是上三角阵，经计算不难得到
\[
P^{-1}A^{-1}P=(P^{-1}AP)^{-1}=
\left(\begin{array}{cccc}
\lambda_1^{-1} & * & \cdots & *\\
0 & \lambda_2^{-1} & \cdots & *\\
\vdots & \vdots & & \vdots\\
0 & 0 & \cdots & \lambda_n^{-1}
\end{array}\right).
\]
因此 $A^{-1}$ 的全部特征值为 $\lambda_1^{-1},\lambda_2^{-1},\cdots,\lambda_n^{-1}$. $\square$

\subsection*{习题 6.1}

1. 求下列矩阵的特征值与特征向量：
\[
\begin{array}{ccc}
(1)\ \left(\begin{array}{ccc}
3 & 2 & 2\\
2 & 3 & 2\\
-4 & -4 & -3
\end{array}\right);&
(2)\ \left(\begin{array}{ccc}
2 & -1 & 2\\
5 & -3 & 3\\
-1 & 0 & -2
\end{array}\right);&
(3)\ \left(\begin{array}{ccc}
3 & -1 & -1\\
7 & -2 & -3\\
3 & -1 & -1
\end{array}\right).
\end{array}
\]

2. 设 $\varphi$ 是线性空间 $V$ 上的线性变换，$V$ 有一个直和分解：
\[
V=V_1\oplus V_2\oplus\cdots\oplus V_m,
\]
其中 $V_i$ 都是 $\varphi$-不变子空间.

(1) 设 $\varphi$ 限制在 $V_i$ 上的特征多项式为 $f_i(\lambda)$，求证：$\varphi$ 的特征多项式
\[
f(\lambda)=f_1(\lambda)f_2(\lambda)\cdots f_m(\lambda).
\]

(2) 设 $\lambda_0$ 是 $\varphi$ 的特征值，$V_0=\{v\in V\mid \varphi(v)=\lambda_0v\}$ 为特征子空间，$V_{i,0}=V_i\cap V_0=\{v\in V_i\mid \varphi(v)=\lambda_0v\}$，求证：
\[
V_0=V_{1,0}\oplus V_{2,0}\oplus\cdots\oplus V_{m,0}.
\]

3. 设 $n$ 阶分块对角阵 $A=\operatorname{diag}\{A_1,A_2,\cdots,A_m\}$，其中 $A_i$ 是 $n_i$ 阶矩阵.

(1) 任取 $A_i$ 的特征值 $\lambda_i$ 及其特征向量 $x_i\in\mathbb{C}^{n_i}$，求证：可在 $x_i$ 的上下添加适当多的零，得到非零向量 $\widetilde{x}_i\in\mathbb{C}^n$，使 $A\widetilde{x}_i=\lambda_i\widetilde{x}_i$，即 $\widetilde{x}_i$ 是 $A$ 关于特征值 $\lambda_i$ 的特征向量，称为 $x_i$ 的延拓.

(2) 任取 $A$ 的特征值 $\lambda_0$，并设 $\lambda_0$ 是 $A_{i_1},\cdots,A_{i_r}$ 的特征值，但不是其他 $A_j(1\leq j\leq m,\ j\neq i_1,\cdots,i_r)$ 的特征值，求证：$A$ 关于特征值 $\lambda_0$ 的特征子空间的一组基可取为 $A_{i_k}(1\leq k\leq r)$ 关于特征值 $\lambda_0$ 的特征子空间的一组基的延拓的并集.

4. 请利用第 3 题的结论求下列矩阵的特征值与特征向量：
\[
\begin{array}{ccc}
(1)\ \left(\begin{array}{cccc}
1 & 0 & 0 & 0\\
0 & 2 & 0 & 0\\
0 & 0 & a & b\\
0 & 0 & c & d
\end{array}\right);&
(2)\ \left(\begin{array}{cccc}
0 & 1 & 0 & 0\\
1 & 0 & 0 & 0\\
0 & 0 & 0 & -1\\
0 & 0 & 1 & 0
\end{array}\right);&
(3)\ \left(\begin{array}{cccc}
1 & 3 & 0 & 0\\
4 & 2 & 0 & 0\\
0 & 0 & 3 & 1\\
0 & 0 & 6 & 4
\end{array}\right).
\end{array}
\]

5. 设 $\lambda_1,\lambda_2$ 是矩阵 $A$ 的两个不同的特征值，$\alpha_1,\alpha_2$ 分别是 $\lambda_1,\lambda_2$ 的特征向量，求证：$\alpha_1+\alpha_2$ 必不是 $A$ 的特征向量.

6. 证明：$n$ 阶矩阵 $A$ 以任一 $n$ 维非零列向量为特征向量的充分必要条件是 $A=cI_n$，其中 $c$ 是常数.

7. 设矩阵
\[
A=\left(\begin{array}{ccc}
2 & 1 & 1\\
1 & 2 & 1\\
1 & 1 & a
\end{array}\right)
\]
是可逆阵，其伴随阵 $A^*$ 有一个特征值 $\lambda_0$ 及对应的特征向量 $\alpha=(1,b,1)'$，求 $a,b,\lambda_0$ 的值.

8. 设矩阵
\[
A=\left(\begin{array}{ccc}
a & -1 & c\\
5 & b & 3\\
1-c & 0 & -a
\end{array}\right)
\]
满足 $|A|=-1$，又 $A^*$ 有一个特征值 $\lambda_0$ 且对应的特征向量为 $(-1,-1,1)'$，求 $a,b,c,\lambda_0$ 的值.

9. 求证：

(1) 若 $A$ 是幂零阵，即存在正整数 $k$，使 $A^k=O$，则 $A$ 的特征值全为零；

(2) 若 $A$ 是对合阵，即 $A^2=I_n$，则 $A$ 的特征值只可能为 1 或 $-1$；

(3) 若 $A$ 是幂等阵，即 $A^2=A$，则 $A$ 的特征值只可能为 0 或 1.

10. 设 $\alpha,\beta$ 是两个 $n$ 维非零列向量，满足 $\alpha'\beta=0$，求矩阵 $A=\alpha\beta'$ 的全部特征值.

11. 设 $A$ 是 $n$ 阶矩阵，若 $(A+I_n)^m=O$，求证：$A$ 是可逆阵并求出 $|A|$.

12. 设 $A$ 是 $m\times n$ 矩阵，$B$ 是 $n\times m$ 矩阵，且 $m\geq n$. 求证：
\[
|\lambda I_m-AB|=\lambda^{m-n}|\lambda I_n-BA|.
\]
特别，若 $A,B$ 都是 $n$ 阶矩阵，则 $AB$ 与 $BA$ 有相同的特征多项式.

13. 设 $\alpha,\beta$ 是两个 $n$ 维非零列向量，求矩阵 $I_n-\alpha\beta'$ 的全部特征值.

14. 设 $A,B$ 为 $n$ 阶矩阵，求证：$AB+B,BA+B$ 有相同的特征值.

15. 设 $P$ 是可逆阵，$B=PAP^{-1}-P^{-1}AP$，求证：$B$ 的特征值之和为零.

16. 设 $n$ 阶复矩阵 $A,B$ 乘法可交换，即 $AB=BA$，求证：

(1) $A,B$ 的特征子空间互为不变子空间；

(2) $A,B$ 至少有一个公共的特征向量；

(3) $A,B$ 可以同时上三角化，即存在可逆阵 $P$，使 $P^{-1}AP$ 和 $P^{-1}BP$ 都是上三角阵.

\section*{§6.2\quad 对角化}

我们将在这一节里回答上一节中提出的问题：什么样的矩阵相似于一个对角阵？

从几何上看，若 $n$ 维线性空间 $V$ 上的线性变换 $\varphi$ 在某组基 $\{e_1,e_2,\cdots,e_n\}$ 下的表示矩阵为对角阵：
\begin{equation}
\left(\begin{array}{cccc}
\lambda_1 & & & \\
& \lambda_2 & & \\
& & \ddots & \\
& & & \lambda_n
\end{array}\right),
\tag{6.2.1}
\end{equation}
则称 $\varphi$ 为可对角化线性变换. 此时 $\varphi(e_i)=\lambda_ie_i$，即 $e_1,e_2,\cdots,e_n$ 是 $\varphi$ 的特征向量，于是 $\varphi$ 有 $n$ 个线性无关的特征向量.

反过来，若 $n$ 维线性空间 $V$ 上的线性变换 $\varphi$ 有 $n$ 个线性无关的特征向量 $e_1,e_2,\cdots,e_n$，则这组向量构成了 $V$ 的一组基，且 $\varphi$ 在这组基下的表示矩阵显然是一个对角阵.

这样我们就证明了下述定理.

\noindent\textbf{定理 6.2.1}\quad 设 $\varphi$ 是 $n$ 维线性空间 $V$ 上的线性变换，则 $\varphi$ 可对角化的充分必要条件是 $\varphi$ 有 $n$ 个线性无关的特征向量.

将上述定义和定理转化为代数的语言. 设 $A$ 是 $n$ 阶矩阵，若 $A$ 相似于对角阵，即存在可逆阵 $P$，使 $P^{-1}AP$ 为对角阵，则称 $A$ 为可对角化矩阵.

\noindent\textbf{定理 6.2.2}\quad 设 $A$ 是 $n$ 阶矩阵，则 $A$ 可对角化的充分必要条件是 $A$ 有 $n$ 个线性无关的特征向量.

那么是否任一 $n$ 阶矩阵均有 $n$ 个线性无关的特征向量呢？当然不是！

\noindent\textbf{例 6.2.1}\quad 矩阵
\[
A=\left(\begin{array}{cc}
1 & 1\\
0 & 1
\end{array}\right)
\]
的特征值为 $1,1$. 将 $\lambda=1$ 代入 $(\lambda I_2-A)x=0$，求得 $A$ 的特征向量为
\[
k\left(\begin{array}{c}
1\\
0
\end{array}\right),\quad k\in\mathbb{K}\backslash\{0\}.
\]
这表明 $A$ 只有一个线性无关的特征向量，因此 $A$ 不可对角化.

也可用反证法来证明，若 $A$ 可对角化，由于 $A$ 的特征值是 $1,1$，故 $A$ 将相似于 $I_2$，即存在可逆阵 $P$，使 $P^{-1}AP=I_2$. 于是 $A=PI_2P^{-1}=I_2$，引出矛盾！

现在我们来讨论不同的特征值和它们相应的特征向量有什么关系. 设 $n$ 维线性空间 $V$ 上的线性变换 $\varphi$ 有 $k$ 个不同特征值 $\lambda_1,\lambda_2,\cdots,\lambda_k$，相应的特征子空间为 $V_1,V_2,\cdots,V_k$.

\noindent\textbf{定理 6.2.3}\quad 若 $\lambda_1,\lambda_2,\cdots,\lambda_k$ 为 $n$ 维线性空间 $V$ 上的线性变换 $\varphi$ 的不同的特征值，则
\[
V_1+V_2+\cdots+V_k=V_1\oplus V_2\oplus\cdots\oplus V_k.
\]

\noindent\textbf{证明}\quad 对 $k$ 用数学归纳法. 若 $k=1$，结论显然成立. 现设对 $k-1$ 个不同的特征值 $\lambda_1,\lambda_2,\cdots,\lambda_{k-1}$，它们相应的特征子空间 $V_1,V_2,\cdots,V_{k-1}$ 之和是直和. 我们要证明 $V_1,V_2,\cdots,V_{k-1},V_k$ 之和为直和，这只需证明：
\begin{equation}
V_k\cap (V_1+V_2+\cdots+V_{k-1})=0
\tag{6.2.2}
\end{equation}
即可. 设 $v\in V_k\cap (V_1+V_2+\cdots+V_{k-1})$，则
\begin{equation}
v=v_1+v_2+\cdots+v_{k-1},
\tag{6.2.3}
\end{equation}
其中 $v_i\in V_i(i=1,2,\cdots,k-1)$. 在 (6.2.3) 式两边作用 $\varphi$，得
\[
\varphi(v)=\varphi(v_1)+\varphi(v_2)+\cdots+\varphi(v_{k-1}).
\]
但 $v,v_1,v_2,\cdots,v_{k-1}$ 都是 $\varphi$ 的特征向量或零向量，因此
\begin{equation}
\lambda_kv=\lambda_1v_1+\lambda_2v_2+\cdots+\lambda_{k-1}v_{k-1}.
\tag{6.2.4}
\end{equation}
在 (6.2.3) 式两边乘以 $\lambda_k$ 减去 (6.2.4) 式得
\[
0=(\lambda_k-\lambda_1)v_1+(\lambda_k-\lambda_2)v_2+\cdots+(\lambda_k-\lambda_{k-1})v_{k-1}.
\]
由归纳假设，$V_1+V_2+\cdots+V_{k-1}$ 是直和，因此 $(\lambda_k-\lambda_i)v_i=0$，而 $\lambda_k-\lambda_i\neq 0$，从而 $v_i=0(i=1,2,\cdots,k-1)$. 这就证明了 (6.2.2) 式. $\square$

\noindent\textbf{推论 6.2.1}\quad 线性变换 $\varphi$ 属于不同特征值的特征向量必线性无关.

\noindent\textbf{推论 6.2.2}\quad 若 $n$ 维线性空间 $V$ 上的线性变换 $\varphi$ 有 $n$ 个不同的特征值，则 $\varphi$ 必可对角化.

推论 6.2.2 另外一个等价的说法就是：若线性变换 $\varphi$ 的特征多项式没有重根，则 $\varphi$ 可对角化. 注意推论 6.2.2 只是可对角化的充分条件而非必要条件，比如说纯量变换 $\varphi=cI_V$ 当然可对角化，但 $\varphi$ 的 $n$ 个特征值都是 $c$. 由定理 6.2.3，我们还可以得到可对角化的第二个充分必要条件.

\noindent\textbf{推论 6.2.3}\quad 设 $\varphi$ 是 $n$ 维线性空间 $V$ 上的线性变换，$\lambda_1,\lambda_2,\cdots,\lambda_k$ 是 $\varphi$ 的全部不同的特征值，$V_i(i=1,2,\cdots,k)$ 是特征值 $\lambda_i$ 的特征子空间，则 $\varphi$ 可对角化的充分必要条件是
\[
V=V_1\oplus V_2\oplus\cdots\oplus V_k.
\]

\noindent\textbf{证明}\quad 先证充分性. 设
\[
V=V_1\oplus V_2\oplus\cdots\oplus V_k,
\]
分别取 $V_i$ 的一组基 $\{e_{i1},e_{i2},\cdots,e_{it_i}\}(i=1,2,\cdots,k)$，则由定理 3.9.3 知这些向量拼成了 $V$ 的一组基，并且它们都是 $\varphi$ 的特征向量. 因此 $\varphi$ 有 $n$ 个线性无关的特征向量，从而 $\varphi$ 可对角化.

再证必要性. 设 $\varphi$ 可对角化，则 $\varphi$ 有 $n$ 个线性无关的特征向量 $\{e_1,e_2,\cdots,e_n\}$，它们构成了 $V$ 的一组基. 不失一般性，可设这组基中前 $t_1$ 个是关于特征值 $\lambda_1$ 的特征向量；接下去的 $t_2$ 个是关于特征值 $\lambda_2$ 的特征向量；$\cdots$；最后 $t_k$ 个是关于特征值 $\lambda_k$ 的特征向量. 对任一 $\alpha\in V$，设 $\alpha=a_1e_1+a_2e_2+\cdots+a_ne_n$，则 $\alpha$ 可写成 $V_1,V_2,\cdots,V_k$ 中向量之和，因此由定理 6.2.3 可得
\[
V=V_1+V_2+\cdots+V_k=V_1\oplus V_2\oplus\cdots\oplus V_k.
\]
$\square$

为了易于从计算的层面判定可对角化，我们引入特征值的度数和重数的概念.

\noindent\textbf{定义 6.2.1}\quad 设 $\varphi$ 是 $n$ 维线性空间 $V$ 上的线性变换，$\lambda_0$ 是 $\varphi$ 的一个特征值，$V_0$ 是属于 $\lambda_0$ 的特征子空间，称 $\dim V_0$ 为 $\lambda_0$ 的度数或几何重数. $\lambda_0$ 作为 $\varphi$ 的特征多项式根的重数称为 $\lambda_0$ 的重数或代数重数.

由线性映射的维数公式可知，特征值 $\lambda_0$ 的度数 $\dim V_0=\dim\operatorname{Ker}(\lambda_0I_V-\varphi)=n-r(\lambda_0I_V-\varphi)$，而特征值 $\lambda_0$ 的重数则由特征多项式 $|\lambda I_V-\varphi|$ 的因式分解决定. 一般来说，特征值的度数和重数之间有如下的不等式关系.

\noindent\textbf{引理 6.2.1}\quad 设 $\varphi$ 是 $n$ 维线性空间 $V$ 上的线性变换，$\lambda_0$ 是 $\varphi$ 的一个特征值，则 $\lambda_0$ 的度数总是小于等于 $\lambda_0$ 的重数.

\noindent\textbf{证明}\quad 设特征值 $\lambda_0$ 的重数为 $m$，度数为 $t$，又 $V_0$ 是属于 $\lambda_0$ 的特征子空间，则 $\dim V_0=t$. 设 $\{e_1,\cdots,e_t\}$ 是 $V_0$ 的一组基. 由于 $V_0$ 中的非零向量都是 $\varphi$ 关于 $\lambda_0$ 的特征向量，故
\[
\varphi(e_i)=\lambda_0e_i,\quad i=1,\cdots,t.
\]
将 $\{e_1,\cdots,e_t\}$ 扩充为 $V$ 的一组基，记为 $\{e_1,\cdots,e_t,e_{t+1},\cdots,e_n\}$，则 $\varphi$ 在这组基下的表示矩阵为
\begin{equation}
A=\left(\begin{array}{cc}
\lambda_0I_t & *\\
O & B
\end{array}\right),
\tag{6.2.5}
\end{equation}
其中 $B$ 是一个 $n-t$ 阶方阵. 因此，线性变换 $\varphi$ 的特征多项式具有如下形状：
\[
|\lambda I_V-\varphi|=|\lambda I_n-A|=(\lambda-\lambda_0)^t|\lambda I_{n-t}-B|,
\]
这表明 $\lambda_0$ 的重数至少为 $t$，即 $t\leq m$. $\square$

\noindent\textbf{定义 6.2.2}\quad 设 $\varphi$ 是 $n$ 维线性空间 $V$ 上的线性变换，若 $\varphi$ 的任一特征值的度数等于重数，则称 $\varphi$ 有完全的特征向量系.

下面我们给出可对角化的第三个充分必要条件.

\noindent\textbf{定理 6.2.4}\quad 设 $\varphi$ 是 $n$ 维线性空间 $V$ 上的线性变换，则 $\varphi$ 可对角化的充分必要条件是 $\varphi$ 有完全的特征向量系.

\noindent\textbf{证明}\quad 设 $\lambda_1,\lambda_2,\cdots,\lambda_k$ 是 $\varphi$ 的全部不同的特征值，它们对应的特征子空间、重数和度数分别记为 $V_i,m_i,t_i(i=1,2,\cdots,k)$. 由重数的定义以及引理 6.2.1 可知
\[
m_1+m_2+\cdots+m_k=n,\quad t_i\leq m_i,\quad i=1,2,\cdots,k.
\]
由推论 6.2.3，我们只要证明 $\varphi$ 有完全的特征向量系当且仅当 $V=V_1\oplus V_2\oplus\cdots\oplus V_k$.

若 $V=V_1\oplus V_2\oplus\cdots\oplus V_k$，则
\[
\begin{aligned}
n&=\dim V=\dim(V_1\oplus V_2\oplus\cdots\oplus V_k)\\
&=\dim V_1+\dim V_2+\cdots+\dim V_k\\
&=\sum_{i=1}^k t_i\leq \sum_{i=1}^k m_i=n,
\end{aligned}
\]
因此 $t_i=m_i(i=1,2,\cdots,k)$，即 $\varphi$ 有完全的特征向量系. 反过来，若 $\varphi$ 有完全的特征向量系，则
\[
\dim(V_1\oplus V_2\oplus\cdots\oplus V_k)=\sum_{i=1}^k t_i=\sum_{i=1}^k m_i=n=\dim V,
\]
从而 $V=V_1\oplus V_2\oplus\cdots\oplus V_k$ 成立. $\square$

\noindent\textbf{注}\quad 在本节中，除了可对角化的第一个充分必要条件是对矩阵和线性变换都进行了阐述之外，其他的定义和结果只是对线性变换进行了阐述和证明，它们对应的关于矩阵的代数版本都可以类似地给出，我们把这些工作留给读者去完成.

例 6.2.1 中的矩阵有一个二重特征值 1，但只有一个线性无关的特征向量. 因此它没有完全的特征向量系，从而它不可能与一个对角阵相似.

已知可对角化矩阵 $A$，如何求出 $P$ 使 $P^{-1}AP$ 是对角阵？下面我们来讨论这个问题. 设 $A$ 的特征值为 $\lambda_1,\lambda_2,\cdots,\lambda_n$，可逆阵 $P=(\alpha_1,\alpha_2,\cdots,\alpha_n)$ 为其列分块. 因为
\[
P^{-1}AP=\operatorname{diag}\{\lambda_1,\lambda_2,\cdots,\lambda_n\},
\]
所以
\[
AP=P\operatorname{diag}\{\lambda_1,\lambda_2,\cdots,\lambda_n\},
\]
即
\[
(A\alpha_1,A\alpha_2,\cdots,A\alpha_n)=(\lambda_1\alpha_1,\lambda_2\alpha_2,\cdots,\lambda_n\alpha_n).
\]
因此 $A\alpha_i=\lambda_i\alpha_i$，即 $\alpha_i$ 就是属于特征值 $\lambda_i$ 的特征向量，于是 $P$ 的 $n$ 个列向量就是 $A$ 的 $n$ 个线性无关的特征向量. 这表明，只要我们求出 $A$ 的 $n$ 个线性无关的特征向量，并将它们按列分块的方式拼成一个矩阵就是要求的 $P$.

\noindent\textbf{注}\quad 因为特征向量不唯一，所以 $P$ 不唯一. 另外，还要注意 $P$ 的第 $i$ 个列向量对应于 $A$ 的第 $i$ 个特征值.

\noindent\textbf{例 6.2.2}\quad 判断矩阵 $A$ 是否相似于对角阵，如是，求出可逆阵 $P$，使 $P^{-1}AP$ 为对角阵：
\[
A=\left(\begin{array}{ccc}
1 & 0 & 0\\
-2 & 5 & -2\\
-2 & 4 & -1
\end{array}\right).
\]

\noindent\textbf{解}\quad 先计算 $A$ 的特征值：
\[
|\lambda I_3-A|=
\left|\begin{array}{ccc}
\lambda-1 & 0 & 0\\
2 & \lambda-5 & 2\\
2 & -4 & \lambda+1
\end{array}\right|
=(\lambda-1)^2(\lambda-3),
\]
$A$ 有特征值 1 (二重) 及 3 (一重). 当 $\lambda=1$ 时，$(\lambda I_3-A)x=0$ 为
\[
\left\{\begin{array}{l}
2x_1-4x_2+2x_3=0,\\
2x_1-4x_2+2x_3=0.
\end{array}\right.
\]
显然这个方程组的系数矩阵秩为 1，因此解空间维数等于 2. 不难求得方程组的基础解系为
\[
\beta_1=\left(\begin{array}{c}
2\\
1\\
0
\end{array}\right),\quad
\beta_2=\left(\begin{array}{c}
-1\\
0\\
1
\end{array}\right).
\]
当 $\lambda=3$ 时，不难求得方程组 $(3I_3-A)x=0$ 的基础解系为 (只有一个向量)
\[
\beta_3=\left(\begin{array}{c}
0\\
1\\
1
\end{array}\right).
\]
因此
\[
P=\left(\begin{array}{ccc}
2 & -1 & 0\\
1 & 0 & 1\\
0 & 1 & 1
\end{array}\right),\quad
P^{-1}AP=\left(\begin{array}{ccc}
1 & 0 & 0\\
0 & 1 & 0\\
0 & 0 & 3
\end{array}\right).
\]

\noindent\textbf{例 6.2.3}\quad 计算 $A^{10}$.
\[
A=\left(\begin{array}{cc}
1 & 0\\
1 & -2
\end{array}\right).
\]

\noindent\textbf{解}\quad 用上例的方法求得
\[
P=\left(\begin{array}{cc}
3 & 0\\
1 & 1
\end{array}\right),\quad
P^{-1}AP=\left(\begin{array}{cc}
1 & 0\\
0 & -2
\end{array}\right).
\]
因此
\[
A^{10}=P\left(\begin{array}{cc}
1 & 0\\
0 & -2
\end{array}\right)^{10}P^{-1}
=\left(\begin{array}{cc}
1 & 0\\
\dfrac{1}{3}(1-2^{10}) & 2^{10}
\end{array}\right).
\]

\subsection*{习题 6.2}

1. 判断下列矩阵能否相似于对角阵，并说明理由：
\[
\begin{array}{ccc}
(1)\ \left(\begin{array}{ccc}
3 & 2 & 2\\
2 & 3 & 2\\
-4 & -4 & -3
\end{array}\right);&
(2)\ \left(\begin{array}{ccc}
2 & -1 & 2\\
5 & -3 & 3\\
-1 & 0 & -2
\end{array}\right);&
(3)\ \left(\begin{array}{ccc}
3 & -1 & -1\\
7 & -2 & -3\\
3 & -1 & -1
\end{array}\right).
\end{array}
\]

2. 求可逆阵 $P$，使 $P^{-1}AP$ 是对角阵：
\[
\begin{array}{ccc}
(1)\ A=\left(\begin{array}{ccc}
-1 & 4 & -2\\
-3 & 4 & 0\\
-3 & 1 & 3
\end{array}\right);&
(2)\ A=\left(\begin{array}{ccc}
1 & -2 & 2\\
2 & -4 & 4\\
1 & -2 & 2
\end{array}\right);&
(3)\ A=\left(\begin{array}{ccc}
0 & 1 & 0\\
1 & 0 & 0\\
0 & 0 & 1
\end{array}\right).
\end{array}
\]

3. 设三阶矩阵 $A$ 的特征值为 $1,1,4$，对应的特征向量依次为 $(2,1,0)'$，$(-1,0,1)'$，$(0,1,1)'$，试求 $A$.

4. 设 $A$ 是三阶矩阵，$A\alpha_j=j\alpha_j(j=1,2,3)$，其中 $\alpha_1=(1,2,1)'$；$\alpha_2=(0,-2,1)'$；$\alpha_3=(1,1,2)'$，试求 $A$.

5. 已知矩阵
\[
A=\left(\begin{array}{ccc}
1 & -1 & 1\\
2 & x & -2\\
-3 & -3 & y
\end{array}\right),\quad
B=\left(\begin{array}{ccc}
2 & 0 & 0\\
0 & 2 & 0\\
0 & 0 & z
\end{array}\right)
\]
相似.

(1) 求 $x,y,z$ 的值；

(2) 求一个满足 $P^{-1}AP=B$ 的可逆阵 $P$.

6. 设
\[
A=\left(\begin{array}{ccc}
3 & 2 & -2\\
-k & -1 & k\\
4 & 2 & -3
\end{array}\right),
\]
当 $k$ 为何值时，存在可逆阵 $P$，使 $P^{-1}AP$ 是对角阵？求出 $P$ 和对角阵.

7. 设矩阵
\[
A=\left(\begin{array}{ccc}
1 & -1 & 1\\
2 & 4 & -2\\
-3 & -3 & 5
\end{array}\right),
\]
求 $A^n$.

8. 设 $A=\operatorname{diag}\{A_1,A_2,\cdots,A_m\}$ 是分块对角阵，其中 $A_i$ 都是方阵，求证：$\operatorname{diag}\{A_1,A_2,\cdots,A_m\}$ 相似于 $\operatorname{diag}\{A_{i_1},A_{i_2},\cdots,A_{i_m}\}$，其中 $A_{i_1},A_{i_2},\cdots,A_{i_m}$ 是 $A_1,A_2,\cdots,A_m$ 的一个排列.

9. 设 $A,B$ 为 $n$ 阶方阵，求证：$AB-BA$ 必不相似于 $kI_n$，其中 $k$ 是非零常数.

10. 设 $A$ 是实二阶矩阵且 $|A|<0$，求证：$A$ 实相似于对角阵.

11. 设 $n$ 阶矩阵 $A,B$ 有相同的特征值，且这 $n$ 个特征值互不相等. 求证：存在 $n$ 阶矩阵 $P,Q$，使 $A=PQ,B=QP$.

12. 求证：

(1) 若 $n$ 阶矩阵 $A$ 适合 $A^2=I_n$，则 $A$ 必可对角化；

(2) 若 $n$ 阶矩阵 $A$ 适合 $A^2=A$，则 $A$ 必可对角化.

13. 若矩阵 $A,B$ 有完全的特征向量系，求证：
\[
\left(\begin{array}{cc}
A & O\\
O & B
\end{array}\right)
\]
也有完全的特征向量系.

14. 设 $n$ 阶矩阵
\[
A=\left(\begin{array}{cc}
I_r & B\\
O & -I_{n-r}
\end{array}\right),
\]
求证：$A$ 可对角化.

15. 求证：

(1) 若 $n$ 阶矩阵 $A$ 的特征值都是 $\lambda_0$，但 $A$ 不是纯量阵，则 $A$ 不可对角化. 特别地，非零的幂零阵不可对角化.

(2) 若 $n$ 阶实矩阵 $A$ 适合 $A^2+A+I_n=O$，则 $A$ 在实数域上不可对角化.

16. 设 $n(n>1)$ 阶矩阵 $A$ 的秩为 1，求证：$A$ 可对角化的充分必要条件是 $\operatorname{tr}(A)\neq 0$.

\section*{§6.3\quad 极小多项式与 Cayley-Hamilton 定理}

我们已经知道，数域 $\mathbb{K}$ 上的 $n$ 阶矩阵全体组成了 $\mathbb{K}$ 上的线性空间，其维数等于 $n^2$. 因此对任一 $n$ 阶矩阵 $A$，下列 $n^2+1$ 个矩阵必线性相关：
\[
A^{n^2},A^{n^2-1},\cdots,A,I_n.
\]
也就是说，存在 $\mathbb{K}$ 中不全为零的数 $c_i(i=0,1,2,\cdots,n^2)$，使
\[
c_{n^2}A^{n^2}+c_{n^2-1}A^{n^2-1}+\cdots+c_1A+c_0I_n=O.
\]
这表明矩阵 $A$ 适合数域 $\mathbb{K}$ 上的一个非零多项式.

\noindent\textbf{定义 6.3.1}\quad 若 $n$ 阶矩阵 $A$ (或 $n$ 维线性空间 $V$ 上的线性变换 $\varphi$) 适合一个非零首一多项式 $m(x)$，且 $m(x)$ 是 $A$ (或 $\varphi$) 所适合的非零多项式中次数最小者，则称 $m(x)$ 是 $A$ (或 $\varphi$) 的一个极小多项式或最小多项式.

从本节开始的说明我们知道，极小多项式肯定是存在的，那它唯一吗？

\noindent\textbf{引理 6.3.1}\quad 若 $f(x)$ 是 $A$ 适合的一个多项式，则 $A$ 的极小多项式 $m(x)$ 整除 $f(x)$.

\noindent\textbf{证明}\quad 由多项式的带余除法知道
\[
f(x)=m(x)q(x)+r(x),
\]
且 $\deg r(x)<\deg m(x)$. 将 $x=A$ 代入上式得 $r(A)=O$，若 $r(x)\neq 0$，则 $A$ 适合一个比 $m(x)$ 次数更小的非零多项式，矛盾. 故 $r(x)=0$，即 $m(x)\mid f(x)$. $\square$

\noindent\textbf{命题 6.3.1}\quad 任一 $n$ 阶矩阵的极小多项式必唯一.

\noindent\textbf{证明}\quad 若 $m(x),g(x)$ 都是矩阵 $A$ 的极小多项式，则由上述引理知道 $m(x)$ 能够整除 $g(x)$，$g(x)$ 也能够整除 $m(x)$. 因此 $m(x)$ 与 $g(x)$ 只差一个常数因子，又极小多项式必须首项系数为 1，故 $g(x)=m(x)$. $\square$

\noindent\textbf{例 6.3.1}\quad (1) 纯量阵 $A=cI_n$ 的极小多项式 $m(x)=x-c$.

(2) 方阵 $A=\left(\begin{array}{cc}0&1\\0&0\end{array}\right)$ 满足 $A^2=O$，因此由引理 6.3.1 知 $A$ 的极小多项式 $m(x)$ 必整除 $x^2$. 因为 $A\neq O$，所以 $m(x)\neq x$，从而只能是 $m(x)=x^2$. 这个例子也告诉我们，矩阵的极小多项式未必是不可约的多项式.

\noindent\textbf{命题 6.3.2}\quad 相似的矩阵具有相同的极小多项式.

\noindent\textbf{证明}\quad 设矩阵 $A$ 和 $B$ 相似，即存在可逆阵 $P$，使 $B=P^{-1}AP$. 设 $A,B$ 的极小多项式分别为 $m(x),g(x)$，注意到
\[
m(B)=m(P^{-1}AP)=P^{-1}m(A)P=O,
\]
因此 $g(x)\mid m(x)$. 同理，$m(x)\mid g(x)$，故 $m(x)=g(x)$. $\square$

\noindent\textbf{命题 6.3.3}\quad 设 $A$ 是一个分块对角阵
\[
A=\left(\begin{array}{cccc}
A_1 & & & \\
& A_2 & & \\
& & \ddots & \\
& & & A_k
\end{array}\right),
\]
其中 $A_i$ 都是方阵，则 $A$ 的极小多项式等于诸 $A_i$ 的极小多项式之最小公倍式.

\noindent\textbf{证明}\quad 设 $A$ 的极小多项式为 $m(x)$，$A_i$ 的极小多项式为 $m_i(x)$，诸 $m_i(x)$ 的最小公倍式为 $g(x)$，则 $g(A_i)=O$，于是
\[
g(A)=\left(\begin{array}{cccc}
g(A_1) & & & \\
& g(A_2) & & \\
& & \ddots & \\
& & & g(A_k)
\end{array}\right)=O,
\]
从而 $m(x)\mid g(x)$. 又因为
\[
m(A)=\left(\begin{array}{cccc}
m(A_1) & & & \\
& m(A_2) & & \\
& & \ddots & \\
& & & m(A_k)
\end{array}\right)=O,
\]
所以对每个 $i$ 有 $m(A_i)=O$，从而 $m_i(x)\mid m(x)$. 又 $g(x)$ 是诸 $m_i(x)$ 的最小公倍式，故 $g(x)\mid m(x)$. 综上所述，$m(x)=g(x)$. $\square$

\noindent\textbf{例 6.3.2}\quad 设 $n$ 阶方阵 $A$ 可对角化，$\lambda_1,\lambda_2,\cdots,\lambda_k$ 是 $A$ 的全部不同的特征值，求 $A$ 的极小多项式.

\noindent\textbf{解}\quad 设 $A$ 的极小多项式为 $m(x)$. 由 $A$ 可对角化知存在可逆阵 $P$，使
\[
P^{-1}AP=B=
\left(\begin{array}{cccc}
B_1 & & & \\
& B_2 & & \\
& & \ddots & \\
& & & B_k
\end{array}\right),
\]
其中 $B_i=\lambda_iI(1\leq i\leq k)$ 为纯量阵. 由例 6.3.1、命题 6.3.2 和命题 6.3.3 可得
\[
m(x)=[x-\lambda_1,x-\lambda_2,\cdots,x-\lambda_k]=(x-\lambda_1)(x-\lambda_2)\cdots(x-\lambda_k).
\]
从上面的例子可以看出，$A$ 的特征值都是极小多项式的根. 事实上，这一结论对任意方阵都是成立的.

\noindent\textbf{引理 6.3.2}\quad 设 $m(x)$ 是 $n$ 阶矩阵 $A$ 的极小多项式，$\lambda_0$ 是 $A$ 的特征值，则
\[
(x-\lambda_0)\mid m(x).
\]

\noindent\textbf{证明}\quad 由 $m(A)=O$ 及命题 6.1.2 可得 $m(\lambda_0)=0$，故结论成立. $\square$

从本节开始的分析知道，$n$ 阶矩阵的极小多项式的次数最多不超过 $n^2$. 但是这个估计实在比较粗，我们可以估计得更精确些.

为了研究一个矩阵可能适合的多项式，我们先看比较简单的情形. 设 $A$ 是一个上三角阵：
\[
A=\left(\begin{array}{ccccc}
\lambda_1 & a_{12} & \cdots & & a_{1n}\\
& \lambda_2 & \cdots & & a_{2n}\\
& & \ddots & & \vdots\\
& & & & \lambda_n
\end{array}\right),
\]
主对角线上的元素 $\lambda_1,\lambda_2,\cdots,\lambda_n$ 正好是 $A$ 的全部特征值. 将 $A$ 依次作用于标准单位列向量 $e_1,e_2,\cdots,e_n$，可得 $n$ 个等式：
\[
Ae_1=\lambda_1e_1,
\]
\[
Ae_2=a_{12}e_1+\lambda_2e_2,
\]
\[
\cdots\cdots\cdots
\]
\[
Ae_i=a_{1i}e_1+\cdots+a_{i-1,i}e_{i-1}+\lambda_ie_i,
\]
\[
\cdots\cdots\cdots
\]
\[
Ae_n=a_{1n}e_1+\cdots+a_{n-1,n}e_{n-1}+\lambda_ne_n.
\]
作
\[
f(x)=(x-\lambda_1)(x-\lambda_2)\cdots(x-\lambda_n),
\]
注意到 $(A-\lambda_iI_n)(A-\lambda_jI_n)=(A-\lambda_jI_n)(A-\lambda_iI_n)$，不难算出：
\[
f(A)e_i=(A-\lambda_1I_n)(A-\lambda_2I_n)\cdots(A-\lambda_nI_n)e_i=0
\]
对一切 $i=1,2,\cdots,n$ 成立，因此 $f(A)=O$. 注意到 $f(x)$ 是 $A$ 的特征多项式，因此 $A$ 适合它的特征多项式. 利用定理 6.1.2，我们很容易把上述结论推广到一般的情形.

\noindent\textbf{定理 6.3.1 (Cayley-Hamilton 定理)}\quad 设 $A$ 是数域 $\mathbb{K}$ 上的 $n$ 阶矩阵，$f(x)$ 是 $A$ 的特征多项式，则 $f(A)=O$.

\noindent\textbf{证明}\quad 由定理 6.1.2 知 $A$ 复相似于一个上三角阵，也就是说存在可逆阵 $P$，使 $P^{-1}AP=B$ 是一个上三角阵，其中 $P$ 与 $B$ 都是复矩阵，但 $A$ 与 $B$ 有相同的特征多项式 $f(x)$. 记
\[
f(x)=x^n+a_1x^{n-1}+\cdots+a_n,
\]
则 $f(B)=O$. 而
\[
\begin{aligned}
f(A)&=A^n+a_1A^{n-1}+\cdots+a_nI_n\\
&=(PBP^{-1})^n+a_1(PBP^{-1})^{n-1}+\cdots+a_nI_n\\
&=PB^nP^{-1}+a_1PB^{n-1}P^{-1}+\cdots+a_nI_n\\
&=P(B^n+a_1B^{n-1}+\cdots+a_nI_n)P^{-1}\\
&=Pf(B)P^{-1}=O.
\end{aligned}
\]
$\square$

\noindent\textbf{推论 6.3.1}\quad $n$ 阶矩阵 $A$ 的极小多项式是其特征多项式的因式. 特别，$A$ 的极小多项式的次数不超过 $n$.

\noindent\textbf{证明}\quad 由 Cayley-Hamilton 定理及引理 6.3.1 即得结论. $\square$

\noindent\textbf{推论 6.3.2}\quad $n$ 阶矩阵 $A$ 的极小多项式和特征多项式有相同的根 (不计重数).

\noindent\textbf{证明}\quad 由引理 6.3.2 和推论 6.3.1 即得结论. $\square$

虽然说矩阵 $A$ 的极小多项式是其特征多项式的因式，但一般来说并不一定等于特征多项式. 对一个给定的矩阵，如何求它的极小多项式？我们将在下一章中加以阐述.

由于矩阵与线性变换之间有一一对应关系，因此我们有下述推论.

\noindent\textbf{推论 6.3.3 (Cayley-Hamilton 定理)}\quad 设 $\varphi$ 是 $n$ 维线性空间 $V$ 上的线性变换，$f(x)$ 是 $\varphi$ 的特征多项式，则 $f(\varphi)=0$.

\subsection*{习题 6.3}

1. 设数域 $\mathbb{F}$ 上的 $n$ 阶矩阵 $A$ 的极小多项式为 $m(x)$，求证：$\mathbb{F}[A]=\{f(A)\mid f(x)\in\mathbb{F}[x]\}$ 是 $M_n(\mathbb{F})$ 的子空间，且 $\dim\mathbb{F}[A]=\deg m(x)$.

2. 求证：任一矩阵与其转置均有相同的极小多项式.

3. 举例说明：极小多项式和特征多项式都相同的矩阵未必相似.

4. 设 $m(x)$ 和 $f(x)$ 分别是 $n$ 阶矩阵 $A$ 的极小多项式和特征多项式，求证：$f(x)\mid m(x)^n$.

5. 设 $n(n>1)$ 阶矩阵 $A$ 的秩为 1，求证：$A$ 的极小多项式为 $x^2-\operatorname{tr}(A)x$.

6. 设 $f(x)$ 和 $m(x)$ 分别是 $m$ 阶矩阵 $A$ 的特征多项式和极小多项式，$g(x)$ 和 $n(x)$ 分别是 $n$ 阶矩阵 $B$ 的特征多项式和极小多项式，证明以下结论等价：

(1) $A,B$ 没有公共的特征值；

(2) $(f(x),g(x))=1$ 或 $(f(x),n(x))=1$ 或 $(m(x),g(x))=1$ 或 $(m(x),n(x))=1$；

(3) $f(B)$ 或 $m(B)$ 或 $g(A)$ 或 $n(A)$ 是可逆阵.

7. 设 $f(x)$ 和 $m(x)$ 分别是 $n$ 阶矩阵 $A$ 的特征多项式和极小多项式，$g(x)$ 是一个多项式，求证：$g(A)$ 是可逆阵的充分必要条件是 $(f(x),g(x))=1$ 或 $(m(x),g(x))=1$.

8. 证明：$n$ 阶方阵 $A$ 为可逆阵的充分必要条件是 $A$ 的极小多项式的常数项不为零.

9. 设 $A$ 是 $n$ 阶可逆阵，求证：$A^{-1}=g(A)$，其中 $g(x)$ 是一个 $n-1$ 次多项式.

10. 设 $A$ 为 $m$ 阶矩阵，$B$ 为 $n$ 阶矩阵，求证：若 $A,B$ 没有公共的特征值，则矩阵方程
\[
AX=XB
\]
只有零解 $X=O$.

11. 设 $A,B$ 分别为 $m,n$ 阶矩阵，$V$ 为 $m\times n$ 矩阵全体构成的线性空间，$V$ 上的线性变换 $\varphi$ 定义为：$\varphi(X)=AX-XB$. 求证：$\varphi$ 是线性自同构的充分必要条件是 $A,B$ 没有公共的特征值. 此时，对任一 $m\times n$ 矩阵 $C$，矩阵方程 $AX-XB=C$ 存在唯一解.

12. 设 $\varphi$ 是数域 $\mathbb{K}$ 上 $n$ 维线性空间 $V$ 上的线性变换，其特征多项式是 $f(\lambda)$ 且 $f(\lambda)=f_1(\lambda)f_2(\lambda)$，其中 $f_1(\lambda),f_2(\lambda)$ 是互素的首一多项式. 令 $V_1=\operatorname{Ker}f_1(\varphi),V_2=\operatorname{Ker}f_2(\varphi)$，求证：

(1) $V_1,V_2$ 是 $\varphi$-不变子空间且 $V=V_1\oplus V_2$；

(2) $V_1=\operatorname{Im}f_2(\varphi),V_2=\operatorname{Im}f_1(\varphi)$；

(3) $\varphi|_{V_1}$ 的特征多项式是 $f_1(\lambda)$，$\varphi|_{V_2}$ 的特征多项式是 $f_2(\lambda)$.

\section*{* §6.4\quad 特征值的估计}

在许多实际问题及理论问题中，常常需要对矩阵的特征值做出估计. 比如，特征值是否在单位圆内？特征值的实部是否小于零？我们将在这一节中介绍两个常用的定理.

一般来说，矩阵的特征值是一些复数. 复数值的估计常常用复平面上的圆来给定范围. 复平面上以 $z_0$ 为圆心、$r$ 为半径的圆常用 $|z-z_0|=r$ 来表示，该圆内部 (包括圆周) 用 $|z-z_0|\leq r$ 来表示，该圆的外部用 $|z-z_0|>r$ 来表示.

现设 $A=(a_{ij})_{n\times n}$ 是一个 $n$ 阶方阵，$A$ 的特征多项式为
\[
f(\lambda)=|\lambda I_n-A|.
\]
令
\[
R_i=\sum_{\substack{j=1\\j\neq i}}^n |a_{ij}|=|a_{i1}|+\cdots+|a_{i,i-1}|+|a_{i,i+1}|+\cdots+|a_{in}|,
\]
即 $R_i$ 为 $A$ 的第 $i$ 行元素去掉 $a_{ii}$ 后的模长之和. 我们有下列“圆盘定理”(又称 Gerschgorin (戈尔斯哥利) 圆盘第一定理).

\noindent\textbf{定理 6.4.1}\quad 设 $A=(a_{ij})_{n\times n}$ 是 $n$ 阶复矩阵，则 $A$ 的特征值在复平面上的下列圆盘 (又称戈氏圆盘) 中：
\[
|z-a_{ii}|\leq R_i,\quad i=1,2,\cdots,n.
\]

\noindent\textbf{证明}\quad 任取 $A$ 的一个特征值 $\lambda_0$，设 $\xi$ 为属于 $\lambda_0$ 的特征向量，则 $A\xi=\lambda_0\xi$. 记 $\xi$ 的第 $i$ 个坐标为 $x_i(i=1,2,\cdots,n)$，将 $A\xi=\lambda_0\xi$ 写成线性方程组：
\[
\left\{\begin{array}{l}
a_{11}x_1+a_{12}x_2+\cdots+a_{1n}x_n=\lambda_0x_1,\\
a_{21}x_1+a_{22}x_2+\cdots+a_{2n}x_n=\lambda_0x_2,\\
\qquad\qquad\cdots\cdots\cdots\\
a_{n1}x_1+a_{n2}x_2+\cdots+a_{nn}x_n=\lambda_0x_n.
\end{array}\right.
\]
设 $|x_1|,|x_2|,\cdots,|x_n|$ 中 $|x_r|$ 最大，从上式可得
\[
(\lambda_0-a_{rr})x_r=a_{r1}x_1+\cdots+a_{r,r-1}x_{r-1}+a_{r,r+1}x_{r+1}+\cdots+a_{rn}x_n.
\]
于是
\[
\begin{aligned}
|\lambda_0-a_{rr}||x_r|
&\leq |a_{r1}||x_1|+\cdots+|a_{r,r-1}||x_{r-1}|+|a_{r,r+1}||x_{r+1}|+\cdots+|a_{rn}||x_n|\\
&\leq (|a_{r1}|+\cdots+|a_{r,r-1}|+|a_{r,r+1}|+\cdots+|a_{rn}|)|x_r|.
\end{aligned}
\]
此即
\[
|\lambda_0-a_{rr}||x_r|\leq R_r|x_r|.
\]
由于 $|x_r|>0$，故
\[
|\lambda_0-a_{rr}|\leq R_r.
\]
$\square$

\noindent\textbf{例 6.4.1}\quad 估计下列矩阵特征值的范围：
\[
\left(\begin{array}{rrrr}
1 & 0.5 & -0.2 & -1\\
0.3 & 2 & -0.2 & 1.1\\
-0.5 & 0.1 & -4 & 0.2\\
-1 & -0.1 & 0.2 & 0
\end{array}\right).
\]

\noindent\textbf{解}\quad 由定理 6.4.1，写出 4 个戈氏圆盘为
\[
\begin{array}{ll}
D_1: & |z-1|\leq 0.5+0.2+1=1.7,\\
D_2: & |z-2|\leq 0.3+0.2+1.1=1.6,\\
D_3: & |z+4|\leq 0.5+0.1+0.2=0.8,\\
D_4: & |z|\leq 1+0.1+0.2=1.3.
\end{array}
\]
若把这 4 个圆盘画在复平面上，则可看出 $D_1,D_2,D_4$ 连在一起，$D_3$ 不与其他 3 个圆盘相连.

若一个戈氏圆盘与另一个相连，则称这两个圆盘内 (包括圆周) 的区域是连通的. 若几个圆盘连在一起，比如 $D_1$ 与 $D_2$ 相连，$D_2$ 与 $D_4$ 相连，则称这些相连圆盘内的区域 (包括圆周) 为连通区域，这几个圆盘称为连通圆盘.

在阐述戈氏圆盘第二定理之前，我们需要引用如下的结果.

\noindent\textbf{定理 6.4.2}\quad 设
\[
f(x)=x^n+a_1x^{n-1}+\cdots+a_{n-1}x+a_n
\]
是 $n$ 次首一复系数多项式，则 $f(x)$ 的 $n$ 个根 $\lambda_1,\lambda_2,\cdots,\lambda_n$ 作为整体是 $a_1,a_2,\cdots,a_n$ 的连续函数.

我们先解释一下定理的含义. 设 $f(x)$ 的 $n$ 个根为 $\lambda_1,\lambda_2,\cdots,\lambda_n$，它们都是 $a_1,a_2,\cdots,a_n$ 的函数. $\lambda_1,\lambda_2,\cdots,\lambda_n$ 作为整体关于 $a_1,a_2,\cdots,a_n$ 连续的意思是：对任意给定的 $\varepsilon>0$，存在 $\delta>0$，使得对满足条件 $|\widetilde{a}_i-a_i|<\delta(i=1,2,\cdots,n)$ 的任意多项式
\[
\widetilde{f}(x)=x^n+\widetilde{a}_1x^{n-1}+\cdots+\widetilde{a}_{n-1}x+\widetilde{a}_n,
\]
$\widetilde{f}(x)$ 的 $n$ 个根存在一个排列，记为 $\widetilde{\lambda}_1,\widetilde{\lambda}_2,\cdots,\widetilde{\lambda}_n$，成立
\[
|\widetilde{\lambda}_i-\lambda_i|<\varepsilon,\quad i=1,2,\cdots,n.
\]

关于定理 6.4.2 的证明，有兴趣的读者可以参考 [2]. 注意定理 6.4.2 并非断言多项式的某个根是关于系数的连续函数，因此一般来说，矩阵的某个特征值也并非是关于矩阵元素的连续函数. 例如，矩阵
\[
A(z)=\left(\begin{array}{cc}
0 & z\\
1 & 0
\end{array}\right),
\]
当参数 $z$ 在复平面上的单位圆内变动时，$A(z)$ 的特征值就不是关于 $z$ 的连续函数. 不过当参数的变动范围是实轴上的区间时，我们有如下肯定的结论，其证明可以参考 [3].

\noindent\textbf{定理 6.4.3}\quad 设 $D$ 是实轴上的区间，$A:D\to M_n(\mathbb{C})$ 是一个取值为复矩阵的连续函数，则 $A(t)$ 的 $n$ 个特征值 $\lambda_1(t),\lambda_2(t),\cdots,\lambda_n(t)$ 都是关于 $t$ 的连续函数.

\noindent\textbf{定理 6.4.4}\quad 设矩阵 $A=(a_{ij})_{n\times n}$ 的 $n$ 个戈氏圆盘分成若干个连通区域，若其中一个连通区域含有 $k$ 个戈氏圆盘，则有且只有 $k$ 个特征值落在这个连通区域内 (若两个戈氏圆盘重合，需计重数；又若特征值为重根，也计重数).

\noindent\textbf{证明}\quad 考虑下列带参数 $t$ 的矩阵：
\[
A(t)=\left(\begin{array}{cccc}
a_{11} & ta_{12} & \cdots & ta_{1n}\\
ta_{21} & a_{22} & \cdots & ta_{2n}\\
\vdots & \vdots & & \vdots\\
ta_{n1} & ta_{n2} & \cdots & a_{nn}
\end{array}\right),
\]
显然 $A(1)=A$，而 $A(0)$ 为下列对角阵：
\[
A(0)=\left(\begin{array}{cccc}
a_{11} & & & \\
& a_{22} & & \\
& & \ddots & \\
& & & a_{nn}
\end{array}\right).
\]
由戈氏圆盘第一定理，矩阵 $A(t)$ 的特征值落在下列圆盘中：
\[
|z-a_{ii}|\leq tR_i,\quad i=1,2,\cdots,n,
\]
其中 $R_i$ 为 $A$ 的第 $i$ 行元素去掉 $a_{ii}$ 后的模长之和. 让 $t$ 从 0 变到 1，则 $A(t)$ 的特征值始终不会越出下列圆盘：
\[
|z-a_{ii}|\leq R_i,\quad i=1,2,\cdots,n.
\]
又 $A(t)$ 的特征多项式的系数是 $t$ 的多项式，故由定理 6.4.3 知 $A(t)$ 的特征值关于 $t$ 连续. 若 $k$ 个圆盘组成一个连通区域，由于 $A(0)$ 的 $k$ 个特征值 (即 $a_{ii}$ 中的 $k$ 个元) 总在这 $k$ 个圆盘内，故 $A$ 在这 $k$ 个圆盘内至少有 $k$ 个特征值，即它们不可能跑到与这 $k$ 个圆盘不相连通的圆盘内. 由于这一结论对任一圆盘连通区域都对，故由这 $k$ 个圆盘组成的连通区域内只有 $A$ 的 $k$ 个特征值. $\square$

\subsection*{习题 6.4}

1. 估计下列矩阵的特征值范围并在复平面上画出其戈氏圆盘：
\[
(1)\ \left(\begin{array}{rrr}
1 & -0.1 & 0.2\\
-1.1 & -1 & 0.4\\
-0.3 & 0.1 & 0
\end{array}\right);\qquad
(2)\ \left(\begin{array}{rrrr}
-3 & 1 & 0.1 & -0.5\\
0.1 & 2 & 0.4 & -0.1\\
1 & 0.3 & 0 & 0\\
0 & -2 & 1 & 1
\end{array}\right).
\]

2. 如果 $n$ 阶实方阵 $A=(a_{ij})$ 适合条件：
\[
|a_{ii}|>\sum_{\substack{j=1\\j\neq i}}^n |a_{ij}|,\quad i=1,2,\cdots,n,
\]
则称 $A$ 是严格对角占优阵. 求证：严格对角占优阵的特征值不等于零，从而 $A$ 必是非异阵.

3. 设 $A$ 是主对角元素全部大于零的严格对角占优阵，求证：$A$ 的特征值的实部均为正数，从而 $|A|>0$.

4. 如果圆盘定理中有一个连通分支由两个圆盘外切组成，证明：每个圆盘除去切点的区域不可能同时包含两个特征值.

\subsection*{历史与展望}

特征值和特征向量通常在线性变换或矩阵理论中引入，然而从历史上看，它们首先出现在二次型和微分方程的研究中. 18 世纪，Euler 研究了刚体的旋转运动，发现了主轴的重要性，Lagrange 意识到主轴是惯性矩阵的特征向量. 19 世纪初，Cauchy 发现他们的工作可以用来对二次曲面进行分类，并将其推广到任意维. Cauchy 还创造了“特征根 (characteristic root)”一词，这来源于特征方程.

后来，Fourier 利用 Lagrange 和 Laplace 的工作，在 1822 年出版的著作《查勒尔的理论分析》中通过分离变量来求解热方程. Sturm (斯图姆) 进一步发展了 Fourier 的思想，这引起了 Cauchy 的注意，Cauchy 将这些思想与自己的思想结合起来，得出了实对称阵具有实特征值的事实. 1855 年，Hermite 将其推广到了现在所称的 Hermite 矩阵上. 大约在同一时间，Brioschi (布廖斯基) 证明了正交阵的特征值位于单位圆上，Clebsch (克莱布施) 发现了斜对称阵的相应结果. 与此同时，Liouville 研究了与 Sturm 相似的特征值问题，从他们的工作中发展出来的学科现在被称为 Sturm-Liouville 理论. 19 世纪末，Schwarz 研究了一般区域上 Laplace 方程的第一个特征值.

20 世纪初，Hilbert (希尔伯特) 通过将算子视为无限阶矩阵来研究积分算子的特征值. 1904 年，他第一个使用德语单词 “eigen (意思是自己的)” 来表示特征值和特征向量. 在某一段时间，英文中的标准术语是“固有值 (proper value)”，但今天的标准术语是“特征值 (eigenvalue)”.

在 1858 年的论文《矩阵理论的研究报告》中，Cayley 证明了 Cayley-Hamilton 定理，即方阵适合它自己的特征多项式. 这个证明由他对 $2\times 2$ 矩阵和 $3\times 3$ 矩阵的计算组成. 他指出这个结果可推广到更高阶，但他补充道：“我认为没有必要花这个力气去正式证明这个定理对任意阶矩阵都成立.” Hamilton 在他的四元数研究中独立地证明了这个定理 ($n=4$ 的情形，但是没有使用矩阵的符号). 1878 年 Frobenius (弗罗本纽斯) 给出了 Cayley-Hamilton 定理的一般情形的证明.

圆盘定理由苏联数学家 Gerschgorin 在 1931 年给出. 虽然随着计算机技术的不断发展，现在我们可以利用 MATLAB 等数学软件快捷地计算出矩阵特征值的近似值，但圆盘定理在矩阵特征值的理论估计方面仍然有着用武之地.

\subsection*{复习题六}

1. 设 $V$ 是 $n$ 阶矩阵全体组成的线性空间，$\varphi$ 是 $V$ 上的线性变换：$\varphi(X)=AX$，其中 $A$ 是一个 $n$ 阶矩阵. 求证：$\varphi$ 和 $A$ 具有相同的特征值 (重数可能不同).

2. 设 $A$ 是 $n$ 阶整数矩阵，$p,q$ 为互素的整数且 $q>1$. 求证：矩阵方程
\[
Ax=\frac{p}{q}x
\]
必无非零解.

3. 求下列 $n$ 阶矩阵的特征值：
\[
A=\left(\begin{array}{ccccc}
0 & a & \cdots & a & a\\
b & 0 & \cdots & a & a\\
\vdots & \vdots & & \vdots & \vdots\\
b & b & \cdots & 0 & a\\
b & b & \cdots & b & 0
\end{array}\right).
\]

4. 设 $n$ 阶矩阵 $A$ 的全体特征值为 $\lambda_1,\lambda_2,\cdots,\lambda_n$，求 $2n$ 阶矩阵
\[
\left(\begin{array}{cc}
A & A^2\\
A^2 & A
\end{array}\right)
\]
的全体特征值.

5. 设矩阵 $A$ 的特征多项式为 $f(\lambda)$，$A$ 的全体特征值为 $\lambda_1,\lambda_2,\cdots,\lambda_n$，$g(\lambda)$ 为任一多项式. 证明：矩阵 $g(A)$ 的行列式等于 $f(\lambda),g(\lambda)$ 的结式 $R(f,g)$.

6. 设首一多项式 $f(x)=x^n+a_{n-1}x^{n-1}+\cdots+a_1x+a_0$，$f(x)$ 的友阵
\[
C=\left(\begin{array}{ccccc}
0 & 0 & \cdots & 0 & -a_0\\
1 & 0 & \cdots & 0 & -a_1\\
0 & 1 & \cdots & 0 & -a_2\\
\vdots & \vdots & & \vdots & \vdots\\
0 & 0 & \cdots & 1 & -a_{n-1}
\end{array}\right).
\]

(1) 求证：矩阵 $C$ 的特征多项式就是 $f(\lambda)$.

(2) 设 $f(x)$ 的根为 $\lambda_1,\lambda_2,\cdots,\lambda_n$，$g(x)$ 为任一多项式，求以 $g(\lambda_1),g(\lambda_2),\cdots,g(\lambda_n)$ 为根的 $n$ 次多项式.

7. 求证：$n$ 阶矩阵 $A$ 为幂零阵的充分必要条件是 $A$ 的特征值全为零.

8. 设 $V$ 是数域 $\mathbb{K}$ 上的 $n$ 阶方阵全体构成的线性空间，$n$ 阶方阵
\[
P=\left(\begin{array}{cccc}
0 & \cdots & 0 & 1\\
0 & \cdots & 1 & 0\\
\vdots & & \vdots & \vdots\\
1 & \cdots & 0 & 0
\end{array}\right),
\]
$V$ 上的线性变换 $\eta$ 定义为 $\eta(X)=PX'P$. 试求 $\eta$ 的全体特征值及其特征向量.

9. 设 $n$ 阶方阵 $A$ 的每行每列只有一个元素非零，并且那些非零元素为 1 或 $-1$，证明：$A$ 的特征值都是单位根.

10. 设 $A$ 是 $n$ 阶实方阵，又 $I_n-A$ 的特征值的模长都小于 1，求证：$0<|A|<2^n$.

11. 设 $n$ 阶方阵 $A$ 的全体特征值为 $\lambda_1,\lambda_2,\cdots,\lambda_n$，求证：$A^*$ 的全体特征值为 $\prod_{i\neq 1}\lambda_i,\prod_{i\neq 2}\lambda_i,\cdots,\prod_{i\neq n}\lambda_i$.

12. 设 $\alpha$ 是 $n$ 维实列向量且 $\alpha'\alpha=1$，试求矩阵 $I_n-2\alpha\alpha'$ 的特征值.

13. 设 $A$ 为 $n$ 阶方阵，$\alpha,\beta$ 为 $n$ 维列向量，试求矩阵 $A\alpha\beta'$ 的特征值.

14. 设 $a_i(1\leq i\leq n)$ 都是实数，且 $a_1+a_2+\cdots+a_n=0$，试求下列矩阵的特征值：
\[
A=\left(\begin{array}{cccc}
a_1^2 & a_1a_2+1 & \cdots & a_1a_n+1\\
a_2a_1+1 & a_2^2 & \cdots & a_2a_n+1\\
\vdots & \vdots & & \vdots\\
a_na_1+1 & a_na_2+1 & \cdots & a_n^2
\end{array}\right).
\]

15. 设 $A,B,C$ 分别是 $m\times m,n\times n,m\times n$ 矩阵，满足：$AC=CB,r(C)=r$. 求证：$A$ 和 $B$ 至少有 $r$ 个相同的特征值.

16. 设 $n$ 阶矩阵 $A$ 的特征多项式为
\[
f(\lambda)=\lambda^n+a_1\lambda^{n-1}+\cdots+a_{n-1}\lambda+a_n.
\]
求证：$a_r$ 等于 $(-1)^r$ 乘以 $A$ 的所有 $r$ 阶主子式之和，即
\[
a_r=(-1)^r\sum_{1\leq i_1<i_2<\cdots<i_r\leq n}
A\left(\begin{array}{cccc}
i_1 & i_2 & \cdots & i_r\\
i_1 & i_2 & \cdots & i_r
\end{array}\right),\quad 1\leq r\leq n.
\]
进一步，若设 $A$ 的特征值为 $\lambda_1,\lambda_2,\cdots,\lambda_n$，则
\[
\sum_{1\leq i_1<i_2<\cdots<i_r\leq n}\lambda_{i_1}\lambda_{i_2}\cdots\lambda_{i_r}
=
\sum_{1\leq i_1<i_2<\cdots<i_r\leq n}
A\left(\begin{array}{cccc}
i_1 & i_2 & \cdots & i_r\\
i_1 & i_2 & \cdots & i_r
\end{array}\right),\quad 1\leq r\leq n.
\]

17. 设 $n$ 阶实方阵 $A$ 的特征值全是实数，且 $A$ 的一阶主子式之和与二阶主子式之和都等于零. 求证：$A$ 是幂零阵.

18. 设 $n(n\geq 3)$ 阶非异实方阵 $A$ 的特征值都是实数，且 $A$ 的 $n-1$ 阶主子式之和等于零，证明：存在 $A$ 的一个 $n-2$ 阶主子式，其符号与 $|A|$ 的符号相反.

19. 设 $A$ 是 $n$ 阶对合阵，即 $A^2=I_n$，证明：$n-\operatorname{tr}(A)$ 为偶数，并且 $\operatorname{tr}(A)=n$ 的充分必要条件是 $A=I_n$.

20. 设 4 阶方阵 $A$ 满足：$\operatorname{tr}(A^k)=k(1\leq k\leq 4)$，试求 $A$ 的行列式.

21. 求证：$n$ 阶矩阵 $A$ 是幂零阵的充分必要条件是 $\operatorname{tr}(A^k)=0(1\leq k\leq n)$.

22. 设 $A,B,C$ 是 $n$ 阶矩阵，其中 $C=AB-BA$. 若它们满足条件 $AC=CA,BC=CB$，求证：$C$ 的特征值全为零. 又若将条件减弱为 $ABC=CAB,BAC=CBA$，则上述结论不再成立.

23. 设 $A,B$ 都是 $n$ 阶矩阵且 $AB=BA$. 若 $A$ 是幂零阵，求证：$|A+B|=|B|$.

24. 下列数列称为 Fibonacci (斐波那契) 数列：
\[
a_0=0,\ a_1=1,\ a_2=1,\ a_3=2,\ a_4=3,\ a_5=5,\ a_6=8,\cdots,
\]
通项用递推式来表示为 $a_{n+2}=a_{n+1}+a_n$，试求 Fibonacci 数列通项的显式表达式.

25. 求证：复数域上 $n$ 阶循环矩阵
\[
A=\left(\begin{array}{ccccc}
a_1 & a_2 & a_3 & \cdots & a_n\\
a_n & a_1 & a_2 & \cdots & a_{n-1}\\
\vdots & \vdots & \vdots & & \vdots\\
a_2 & a_3 & a_4 & \cdots & a_1
\end{array}\right)
\]
可对角化，并求出它相似的对角阵及过渡矩阵.

26. 设 $n$ 阶复矩阵 $A$ 可对角化，证明：矩阵
\[
\left(\begin{array}{cc}
A & A^2\\
A^2 & A
\end{array}\right)
\]
也可对角化.

27. 设 $A,B,C$ 都是 $n$ 阶矩阵，$A,B$ 各有 $n$ 个不同的特征值，又 $f(\lambda)$ 是 $A$ 的特征多项式，且 $f(B)$ 是可逆阵. 求证：矩阵
\[
M=\left(\begin{array}{cc}
A & C\\
O & B
\end{array}\right)
\]
相似于对角阵.

28. 设 $A,B$ 是 $n$ 阶矩阵，$A$ 有 $n$ 个不同的特征值，并且 $AB=BA$，求证：$B$ 相似于对角阵.

29. 设 $A,B$ 是 $n$ 阶矩阵，$A$ 有 $n$ 个不同的特征值，并且 $AB=BA$，求证：存在次数不超过 $n-1$ 的多项式 $f(x)$，使 $B=f(A)$.

30. 设 $a,b,c$ 为复数且 $bc\neq 0$，证明下列 $n$ 阶矩阵 $A$ 可对角化：
\[
A=\left(\begin{array}{cccccc}
a & b & & & & \\
c & a & b & & & \\
& c & a & b & & \\
& & \ddots & \ddots & \ddots & \\
& & & c & a & b\\
& & & & c & a
\end{array}\right).
\]

31. 设 $m$ 阶矩阵 $A$ 与 $n$ 阶矩阵 $B$ 没有公共的特征值，且 $A,B$ 均可对角化，又 $C$ 为 $m\times n$ 矩阵，求证：
\[
M=\left(\begin{array}{cc}
A & C\\
O & B
\end{array}\right)
\]
也可对角化.

32. 设 $A$ 为 $m$ 阶矩阵，$B$ 为 $n$ 阶矩阵，$C$ 为 $m\times n$ 矩阵，
\[
M=\left(\begin{array}{cc}
A & C\\
O & B
\end{array}\right),
\]
求证：若 $M$ 可对角化，则 $A,B$ 均可对角化.

33. 设 $A$ 为 $m\times n$ 矩阵，$B$ 为 $n\times m$ 矩阵，又 $|BA|\neq 0$，求证：$AB$ 可对角化的充分必要条件是 $BA$ 可对角化.

34. 设 $A$ 是 $n$ 阶矩阵，求证：伴随阵 $A^*=h(A)$，其中 $h(x)$ 是一个 $n-1$ 次多项式.

35. 设 $n$ 阶方阵 $A,B$ 的特征值全部大于零且满足 $A^2=B^2$，求证：$A=B$.

36. 设 $A=\operatorname{diag}\{A_1,A_2,\cdots,A_m\}$ 为 $n$ 阶分块对角阵，其中 $A_i$ 是 $n_i$ 阶矩阵且两两没有公共的特征值. 设 $B$ 是 $n$ 阶矩阵，满足 $AB=BA$，求证：$B=\operatorname{diag}\{B_1,B_2,\cdots,B_m\}$，其中 $B_i$ 也是 $n_i$ 阶矩阵.

37. 设 $n$ 阶实矩阵 $A$ 的所有特征值都是正实数，证明：对任一实对称阵 $C$，存在唯一的实对称阵 $B$，满足 $A'B+BA=C$.

38. 设 $\varphi$ 是复线性空间 $V$ 上的线性变换，又有两个复系数多项式：
\[
f(x)=x^m+a_1x^{m-1}+\cdots+a_m,\quad g(x)=x^n+b_1x^{n-1}+\cdots+b_n.
\]
设 $\sigma=f(\varphi),\tau=g(\varphi)$，矩阵 $C$ 是 $f(x)$ 的友阵，即
\[
C=\left(\begin{array}{ccccc}
0 & 0 & 0 & \cdots & -a_m\\
1 & 0 & 0 & \cdots & -a_{m-1}\\
0 & 1 & 0 & \cdots & -a_{m-2}\\
\vdots & \vdots & \vdots & & \vdots\\
0 & 0 & 0 & \cdots & -a_1
\end{array}\right).
\]
若 $g(C)$ 是可逆阵，求证：$\operatorname{Ker}\sigma\tau=\operatorname{Ker}\sigma\oplus\operatorname{Ker}\tau$.

\end{document}
